\documentclass[11pt]{article}
\usepackage[a4paper,margin=1.45cm]{geometry}
\usepackage[T1]{fontenc}
\usepackage[utf8]{inputenc}
\usepackage{lmodern}
\usepackage[spanish,es-noquoting,es-noshorthands]{babel}
\usepackage{amsmath,amssymb,mathtools,array,booktabs,pdflscape,adjustbox,diagbox}
\usepackage[protrusion=true,expansion=false]{microtype}
\setlength{\parindent}{1.25em}
\setlength{\parskip}{0pt}
\newcommand{\widebar}[1]{\overline{#1}}
\newcommand{\A}{\Delta}
\newcommand{\D}{\Delta}
\allowdisplaybreaks
\setlength{\emergencystretch}{12em}
\sloppy
\hbadness=10000
\vbadness=10000
\hfuzz=2pt
\vfuzz=10pt
\raggedbottom
\newcommand{\R}{\varrho}
\newcommand{\hata}{\widehat{a}}
\newcommand{\hatv}{\widehat{v}}
\newcommand{\modu}[1]{\;\operatorname{mod} #1}
\newcommand{\mcell}[1]{\ensuremath{\begin{gathered}#1\end{gathered}}}
\newcommand{\pair}[2]{(#1\,|\,#2)}
\newcommand{\eps}{\varepsilon}
\providecommand{\Hgrp}{\mathfrak{H}}
\providecommand{\GaloisRes}{\Phi}
\providecommand{\pol}[2]{\left(#1\frac{\partial}{\partial #2}\right)}
\providecommand{\Deltaop}{\Delta\!\left(\frac{\partial}{\partial x},\frac{\partial}{\partial y},\frac{\partial}{\partial z}\right)}
\providecommand{\Omegaop}{\Omega}
\makeatletter
\renewcommand\footnotesize{\@setfontsize\footnotesize{7.5}{8.7}\selectfont}
\makeatother
\begin{document}

\section*{7. El teorema de finitud para los invariantes de grupos finitos}
\begin{center}
Math. Ann. 77 (1916), pp. 89--92
\end{center}

En lo que sigue se dará una demostración completamente elemental de la finitud de los invariantes de grupos finitos ---una que se apoya sólo en la teoría de las funciones simétricas--- y al mismo tiempo proporciona una especificación efectiva del sistema completo; mientras que la demostración habitual, basada en el teorema de Hilbert sobre bases de módulos (Ann. 36), es sólo una demostración de existencia.\footnote{Cf., por ejemplo, Weber, \emph{Lehrbuch der Algebra} (2.a ed.), vol. 2, \S\ 57.}

Sea el grupo finito $\Hgrp$ compuesto por las $h$ transformaciones lineales (con determinante no nulo) $A_1,\ldots,A_h$, donde $A_k$ denota la transformación lineal
\[
  x_1^{(k)}=\sum_{\nu=1}^{n}a_{1\nu}^{(k)}x_\nu,\ldots,
  x_n^{(k)}=\sum_{\nu=1}^{n}a_{n\nu}^{(k)}x_\nu,
\]
o, abreviadamente, $(x^{(k)})=A_k(x)$. Por tanto, el grupo $\Hgrp$ lleva la fila $(x)$ con elementos $x_1,
\ldots,x_n$ a las filas $(x^{(k)})$ con elementos $x_1^{(k)},
\ldots,x_n^{(k)}$. Puesto que la identidad debe estar contenida entre $A_1,
\ldots,A_h$, la fila $(x)$ también está contenida entre las filas $(x^{(k)})$. Por un invariante racional entero (absoluto) del grupo se entiende una función racional entera de $x_1,
\ldots,x_n$ que permanece idénticamente inalterada bajo la aplicación de $A_1,
\ldots,A_h$; para tal invariante $f(x)$ se tiene, por consiguiente,
\begin{equation}
  f(x)=f(x^{(1)})=\cdots=f(x^{(h)})=
  \frac1h\sum_{k=1}^{h}f(x^{(k)}).
\tag{1}
\end{equation}

\subsection*{1.}
La fórmula (1) dice que $f(x)$ es una función simétrica racional entera de las filas de cantidades $(x^{(1)}),
\ldots,(x^{(h)})$; y, como cada sumando $f(x^{(k)})$ contiene sólo la única fila $(x^{(k)})$, éste es el caso más simple, usualmente llamado el caso de \emph{una forma}. Por el teorema conocido sobre funciones simétricas de filas de cantidades,\footnote{Cf. la nota a 2.} $f(x)$ es, por tanto, representable de manera racional entera mediante las funciones simétricas elementales de esas filas, es decir, mediante los coeficientes $G_{\alpha\alpha_1\ldots\alpha_n}(x)$ de la ``resolvente de Galois''
\[
\begin{aligned}
  \Phi(z,u)
  &=\prod_{k=1}^{h}\bigl(z+u_1x_1^{(k)}+\cdots+u_nx_n^{(k)}\bigr) \\
  &=z^h+
    \sum G_{\alpha\alpha_1\ldots\alpha_n}(x)
    z^{\alpha}u_1^{\alpha_1}\cdots u_n^{\alpha_n},
\end{aligned}
\qquad
\left(\begin{array}{c}
\alpha+\alpha_1+
\cdots+
\alpha_n=h\\
\alpha\ne h
\end{array}\right),
\]
donde los $G_{\alpha\alpha_1\ldots\alpha_n}(x)$ son invariantes de grado $\alpha_1+\cdots+\alpha_n$ en las $x$. Así queda probado:

\emph{Los coeficientes $G_{\alpha\alpha_1\ldots\alpha_n}(x)$ de la resolvente de Galois forman un sistema completo de invariantes del grupo, de tal modo que todo invariante puede representarse racionalmente y enteramente mediante estos invariantes finitamente numerosos.}

\subsection*{2.}
Partiendo de (1), también se puede hacer la siguiente consideración todavía más elemental,\footnote{Esto equivale a una demostración del teorema sobre funciones simétricas de filas de cantidades en el caso de una forma mencionado arriba.} que al mismo tiempo conduce a un segundo sistema completo. Pongamos
\[
  f(x)=a+b x_1^{\mu_1}\cdots x_n^{\mu_n}
       +\cdots+c x_1^{\nu_1}\cdots x_n^{\nu_n},
\]
donde $a,b,\ldots,c$ denotan constantes. Entonces, por (1), se tiene
\[
\begin{aligned}
 h\cdot f(x)
 &=h\cdot a+b\sum_{k=1}^{h}(x_1^{(k)})^{\mu_1}\cdots(x_n^{(k)})^{\mu_n}
     +\cdots
     +c\sum_{k=1}^{h}(x_1^{(k)})^{\nu_1}\cdots(x_n^{(k)})^{\nu_n} \\
 &=h\cdot a+b\cdot J_{\mu_1\ldots\mu_n}+\cdots+c\cdot J_{\nu_1\ldots\nu_n}.
\end{aligned}
\]
Por tanto, todo invariante puede componerse de manera entera y lineal a partir de los invariantes especiales
\[
  J_{\mu_1\ldots\mu_n}
  =\sum_{k=1}^{h}(x_1^{(k)})^{\mu_1}\cdots(x_n^{(k)})^{\mu_n},
\]
y basta probar la afirmación para estos invariantes especiales. Pero, salvo un factor numérico, $J_{\mu_1\ldots\mu_n}$ es el coeficiente de $u_1^{\mu_1}
\cdots u_n^{\mu_n}$ en la expresión
\[
  S_\mu=\sum_{k=1}^{h}
  \bigl(u_1x_1^{(k)}+\cdots+u_nx_n^{(k)}\bigr)^{\mu},
\]
donde $\mu=\mu_1+\cdots+\mu_n$, y esta expresión es la $\mu$-ésima suma de potencias de las $h$ formas lineales
\[
  \xi_1=u_1x_1^{(1)}+\cdots+u_nx_n^{(1)},\ldots,
  \xi_h=u_1x_1^{(h)}+\cdots+u_nx_n^{(h)}.
\]
Se sabe que las infinitamente muchas sumas de potencias $S_\mu$ son combinaciones racionales enteras de
\[
  S_1,S_2,\ldots,S_h,
\]
cuyos coeficientes vienen dados por los invariantes
\[
  J_{\mu_1\ldots\mu_n}\qquad(\mu_1+\cdots+\mu_n\le h).
\]
Así se obtiene un segundo sistema completo:

\emph{Un sistema completo de invariantes del grupo está dado por todos los invariantes $J_{\mu_1\ldots\mu_n}$ para los cuales $\mu_1+\cdots+\mu_n\le h$, donde $h$ denota el orden del grupo.}

La conexión con 1. se establece observando que las sumas de potencias $S_1,
\ldots,S_h$ pueden sustituirse por las funciones simétricas elementales de $\xi_1,
\ldots,
\xi_h$, lo que conduce al sistema completo de los $G_{\alpha\alpha_1\ldots\alpha_n}(x)$ dado allí. Ambos resultados muestran que todos los invariantes pueden expresarse racionalmente y enteramente mediante aquellos cuyo grado en las $x$ no excede el orden $h$ del grupo.

\subsection*{3.}
De lo que se acaba de probar pueden extraerse consecuencias para la representación racional. Todo invariante absoluto racional puede representarse ---como se ve inmediatamente multiplicando numerador y denominador por los conjugados del denominador respecto del grupo--- como cociente de dos invariantes absolutos racionales enteros, no necesariamente primos entre sí. De aquí se sigue que todo invariante absoluto racional puede expresarse racionalmente mediante los coeficientes $G_{\alpha\alpha_1\ldots\alpha_n}(x)$ de la resolvente de Galois $\Phi(z,u)$. Este teorema también puede probarse fácilmente de diversas maneras sin pasar por el teorema de finitud; ya está formulado en Weber II, \S\ 58.\footnote{Sin embargo, la demostración dada allí no es correcta. Sólo se muestra que la función $\Psi(t)$ en la fórmula (7) tiene invariantes como coeficientes, pero no que esos invariantes sean expresables racionalmente mediante los coeficientes de $\Phi(t)$. Esta laguna se evita si, al representar los $x_i^{(k)}$ mediante $\xi_k$, se usa un procedimiento de diferenciación conocido en vez de la fórmula de interpolación de Lagrange. Ésta es la relación obtenida al diferenciar la identidad $\Phi(-\xi_k,u)=0$ con respecto a todos los $u$:
\[
 \left[\frac{\partial\Phi}{\partial u_i}
      -x_i^{(k)}\frac{\partial\Phi}{\partial z}\right]_{z=-\xi_k}=0.
\]
En consecuencia, para cada función racional $\omega(x)$ debe sustituirse a las fórmulas (7) y (8) de Weber la fórmula siguiente:
\[
 \omega(x_1^{(k)},\ldots,x_n^{(k)})=
 \left.
 \omega\!\left(
   \frac{\partial\Phi/\partial u_1}{\partial\Phi/\partial z},\ldots,
   \frac{\partial\Phi/\partial u_n}{\partial\Phi/\partial z}
 \right)\right|_{z=-\xi_k},
\]
la cual contiene sólo coeficientes de $\Phi(z,u)$ y cuya suma sobre todos los $k$, para invariantes $\omega(x)$, da la representación deseada.}

\subsection*{4.}
Finalmente, debe mencionarse que los resultados obtenidos aquí están contenidos implícitamente en mi trabajo ``Cuerpos y sistemas de funciones racionales'';\footnote{Math. Ann. 76, p. 161 (1915).} más precisamente, el sistema completo dado en 1. está contenido en los Teoremas VI y VII, y una demostración de la representación racional independiente de este teorema de finitud está contenida en el Teorema III.\footnote{Para invariantes relativos, los Teoremas VIII y IX contienen una demostración de finitud que igualmente descansa sobre un fundamento distinto del habitual, pero que también es sólo una demostración de existencia; la razón es que los invariantes relativos no forman un cuerpo.} En este caso especial, sin embargo, el curso de la demostración y el resultado se simplifican considerablemente en comparación con la teoría general. Fui llevada a aplicar las investigaciones generales a los invariantes de grupos finitos por conversaciones con E. Fischer, de quien proceden, en sustancia, la demostración dada en 2 ---en el caso general el sistema completo que allí aparece no es definible racionalmente--- y también la nota a 3.

Erlangen, mayo de 1915.

\clearpage
\section*{8. Sobre la representación racional entera de los invariantes de un sistema de arbitrariamente muchas formas fundamentales}
\begin{center}
Math. Ann. 77 (1916), pp. 93--102
\end{center}

Al final de su contribución al volumen de homenaje a Schwarz, ``Sobre los invariantes de un sistema de arbitrariamente muchas formas fundamentales'', D. Hilbert formula la conjetura de que estos invariantes pueden representarse racionalmente y enteramente mediante los finitamente muchos invariantes del sistema $(J,PJ)$, donde $J$ denota el sistema completo de invariantes de las formas fundamentales linealmente independientes, y los invariantes $PJ$ surgen de $J$ mediante procesos polares.

Muestro a continuación, en I, que esta conjetura es de hecho correcta, y ello como una consecuencia simple del teorema de reducción según el cual toda forma en arbitrariamente muchas filas de $n$ variables cada una puede representarse como una suma de polares de formas que contienen sólo $n$ filas fijas y que ellas mismas se derivan de la forma original mediante procesos polares. Este teorema de reducción es un caso especial, formulado por primera vez por F. Mertens, de la generalización, dada por A. Capelli, F. Mertens y J. Deruyts, del desarrollo en serie de Clebsch--Gordan a formas en arbitrariamente muchas filas de $n$ variables cada una; esta generalización se demuestra por transformación formal de los procesos de diferenciación.\footnote{F. Mertens: ``Über eine Formel der Determinantentheorie.'' (Fórmula 5), Sitzb. d. Ak. d. Wiss. Wien, vol. 91, II. Abt. (1885), y más detalladamente (con aplicaciones) en: ``Über ganze Funktionen von $m$ Systemen von je $n$ Unbestimmten''. Monatsh. f. Math. u. Phys. 4.o año (1893). Para algunas de las demostraciones de Capelli (desde 1882), cf., por ejemplo, Math. Ann. 37, o las autografiadas ``Lezioni sulla teoria delle forme algebriche'' (1902). J. Deruyts: \emph{Essai d'une théorie générale des formes algébriques}, Mém. d. 1. Soc. Roy. des Sciences de Liège (2) 17 (1892).}

En II doy además una demostración más conceptual del teorema de reducción, tratando la cuestión como un problema sobre la equivalencia de familias lineales de formas. Este punto de vista, y también una parte esencial de las consideraciones usadas, los tomo de un trabajo ``Sobre sistemas algebraicos de módulos''\footnote{E. Fischer: Über algebraische Modulsysteme und lineare homogene partielle Differentialgleichungen mit konstanten Koeffizienten, \S\ 1--4, J. f. M. 140 (1911).} y de teoremas inéditos de E. Fischer conexos con él, cuyo uso me permitió amablemente.

Finalmente, en III, muestro cómo las consideraciones correspondientes conducen también a los desarrollos en serie más generales.

\subsection*{I.}
Sea $\theta$ una forma homogénea en cada una de las $N>n$ filas
\[
  A_1,A_2,\ldots,A_N,
\]
donde, en general, la fila $A_k$ consta de los $n$ elementos
\[
  A_k^{(1)},A_k^{(2)},\ldots,A_k^{(n)}.
\]
Los coeficientes de $\theta$ pueden ser indeterminadas o cantidades de un dominio de racionalidad dado. Denote además $P$ una función racional entera ---con coeficientes enteros racionales--- de los procesos polares
\[
  P_{hk}=\left(A_h\frac{\partial}{\partial A_k}\right)
  =A_h^{(1)}\frac{\partial}{\partial A_k^{(1)}}+
   A_h^{(2)}\frac{\partial}{\partial A_k^{(2)}}+
   \cdots+
   A_h^{(n)}\frac{\partial}{\partial A_k^{(n)}},
\]
donde el producto $P_{hk}P_{ij}\theta$ debe entenderse como la aplicación de $P_{hk}$ a $P_{ij}\theta$.

Entonces el teorema de reducción mencionado está dado por la identidad
\begin{equation}
  \theta=\sum PZ,
\tag{1}
\end{equation}
donde las formas $Z$ contienen sólo las filas
\[
  A_1,A_2,\ldots,A_n
\]
y se derivan de $\theta$ mediante procesos polares $P$.

Consideremos ahora el sistema de formas fundamentales $(F')$, compuesto por las $N$ formas del mismo orden y con el mismo número de variables
\[
  F_1,F_2,\ldots,F_N,
\]
cuyos coeficientes se toman respectivamente como las $N$ filas
\[
  A_1,A_2,\ldots,A_N.
\]
Sea $(J)$ el sistema completo ---compuesto por finitamente muchos invariantes--- de $F_1,
\ldots,F_n$, de tal modo que todo invariante de $F_1,
\ldots,F_n$ pueda expresarse racionalmente y enteramente mediante los invariantes $(J)$. La conjetura de Hilbert que se ha de probar es entonces que, para los invariantes simultáneos $S$ de $F_1,
\ldots,F_N$, un sistema completo está dado por los finitamente muchos invariantes $(J,PJ)$, donde $PJ$ denota todos los invariantes derivables de $J$ mediante procesos polares $P$.

En efecto, todo invariante simultáneo $S$ con coeficientes enteros racionales de $(F')$ es una forma homogénea en cada una de las $N$ filas $A_1,
\ldots,A_N$ ---con coeficientes enteros racionales--- y por tanto se le puede aplicar la identidad (1). Puesto que se sabe que la aplicación de los procesos polares $P$ a $S$ vuelve a producir invariantes, las formas $Z$ también son invariantes, a saber, invariantes simultáneos de $F_1,
\ldots,F_n$, ya que contienen sólo las filas $A_1,
\ldots,A_n$; de ahí que, por hipótesis, sean funciones racionales enteras de los invariantes $(J)$. Así, la identidad (1) se convierte en
\[
  S=\sum PY,
\]
donde las $Y$ denotan productos de potencias de los invariantes $(J)$. Pero, por la definición de $P$, la ejecución efectiva de los procesos $P$ sobre $Y$ consiste en la ejecución sucesiva, finitamente repetida, de las operaciones polares simples $P_{hk}$, lo que, por la regla de diferenciación de productos, conduce a sumas de productos de invariantes $(J)$ y $(PJ)$. Lo mismo vale para la aplicación de $P_{hk}$ a un producto formado por $(J)$ y $(PJ)$; por tanto, también $PY$ produce sólo combinaciones racionales enteras de $(J,PJ)$. \emph{Con ello queda probada la representabilidad racional entera de todos los invariantes $S$ mediante $(J,PJ)$.}

\subsection*{II.}
Antes de probar el teorema de reducción, recordamos algunos hechos sobre familias lineales de formas. Por una \emph{familia lineal de formas sobre $K$} entendemos un sistema de formas de la misma dimensión, completo en el sentido de que, junto con $\theta_1$ y $\theta_2$, contiene siempre $c_1\theta_1+c_2\theta_2$, donde $c_1,c_2$ son cantidades de un dominio de racionalidad prescrito $K$, y donde $K$ contiene el dominio de coeficientes de las $\theta$. Entre estas formas hay siempre un número finito ---digamos $\rho$--- linealmente independiente sobre $K$, mientras que cualesquiera $\rho+1$ formas satisfacen una relación lineal con coeficientes de $K$; la familia tiene \emph{rango $\rho$} respecto de $K$.\footnote{Se obtienen familias lineales de formas más generales ---cuyo rango no necesita ser finito--- suprimiendo la restricción de igual dimensión, o también suprimiendo la restricción de que $K$ contenga el dominio de coeficientes de las $\theta$.} Usaremos el hecho fácilmente visible de que si $\mathcal L$ y $\mathcal T$ son dos familias lineales sobre $K$ del mismo rango $\rho$, y si $\mathcal T$ es una subfamilia de $\mathcal L$, entonces $\mathcal L$ y $\mathcal T$ son idénticas (equivalentes).

Con esto pasamos al teorema de reducción. La identidad de I,
\[
  \theta=\sum PZ,
\]
pasa a una identidad análoga para $P\theta$ cuando se aplica el mismo proceso polar $P$ a ambos lados; así, el teorema de reducción da al mismo tiempo una representación de todas las polares de $\theta$ por sumas de las polares especiales $PZ$. Como, recíprocamente, estas polares especiales están ciertamente contenidas entre todas las polares, el teorema de reducción afirma la equivalencia de estas dos familias lineales de formas; en particular, también afirma la equivalencia de aquellas subfamilias cuyas formas tienen, en cada fila individual, la misma dimensión que $\theta$. Para probar esta equivalencia introducimos una familia intermedia determinada por las formas $Z$. Por simplicidad damos primero la demostración para el caso de tres filas binarias $(N=3,n=2)$ y luego indicamos brevemente la demostración general, completamente paralela.

Sea entonces $\theta(x,y,z)$ de dimensiones $\alpha,\beta,p$, respectivamente, en las filas
\[
  x_1x_2;\quad y_1y_2;\quad z_1z_2;
\]
y sean primero los coeficientes de $\theta$ indeterminadas $a$ (o números racionales). Consideramos las tres familias lineales de formas siguientes.

1. La familia $\mathcal L$, compuesta por todas las formas $f(x,y,z)$ que surgen de $\theta$ mediante procesos polares $P$ y que, en las filas individuales $x,y,z$, tienen la misma dimensión que $\theta$; en particular, $\mathcal L$ contiene a $\theta$. Si $f(x,y,z)$ se considera como forma en $x,y,z$ y las indeterminadas $a$, entonces el cuerpo de coeficientes es el cuerpo $R$ de los números racionales; también para $K$ debe tomarse $R$, ya que sólo para $\theta_1,\theta_2$ racionales enteras se vuelve $c_1P_1+c_2P_2$ nuevamente un proceso $P$; sea $\mathcal L$ de rango $\rho$ respecto de $R$.

2. La familia $\mathcal S$, compuesta por todas las formas $g(\lambda;x,y)$ definidas por
\[
  g(\lambda;x,y)=f(x,y,\lambda_1x+\lambda_2y),
\]
o, expresadas por procesos polares,
\[
\begin{aligned}
  g(\lambda;x,y)
  &=\frac1{p!}\left[(\lambda_1x+\lambda_2y)\frac{\partial}{\partial z}\right]^p f(x,y,z) \\
  &=g_0(xy)\lambda_1^p+g_1(x,y)\lambda_1^{p-1}\lambda_2+\cdots+g_p(xy)\lambda_2^p,
\end{aligned}
\]
donde
\[
  i!(p-i)!\,g_i(xy)=
  \left(y\frac{\partial}{\partial z}\right)^i
  \left(x\frac{\partial}{\partial z}\right)^{p-i}f(xyz).
\]
\footnote{Como en I,
\[
 \left(t\frac{\partial}{\partial x}\right)=t_1\frac{\partial}{\partial x_1}+t_2\frac{\partial}{\partial x_2},
\]
y por tanto, en particular,
\[
 \left[(\lambda_1x+\lambda_2y)\frac{\partial}{\partial z}\right]
 = (\lambda_1x_1+\lambda_2y_1)\frac{\partial}{\partial z_1}
 + (\lambda_1x_2+\lambda_2y_2)\frac{\partial}{\partial z_2},
\]
\[
 \left[z\left(\frac{\partial}{\partial\lambda_1}\frac{\partial}{\partial x}
 +\frac{\partial}{\partial\lambda_2}\frac{\partial}{\partial y}\right)\right]
 =z_1\left(\frac{\partial}{\partial\lambda_1}\frac{\partial}{\partial x_1}
 +\frac{\partial}{\partial\lambda_2}\frac{\partial}{\partial y_1}\right)
 +z_2\left(\frac{\partial}{\partial\lambda_1}\frac{\partial}{\partial x_2}
 +\frac{\partial}{\partial\lambda_2}\frac{\partial}{\partial y_2}\right).
\]}
Así, los $g_i(xy)$ dan formas $Z$ de la identidad (1). Las $g$ son formas racionales enteras en $x,y,\lambda$ y las indeterminadas $a$; sea $\mathcal S$ de rango $\sigma$ respecto de $R$.

3. La familia $\mathcal T$, obtenida de $\mathcal S$ por el proceso polar inverso del mismo modo que $\mathcal S$ se obtiene de $\mathcal L$; las formas $h$ de $\mathcal T$ están definidas por
\[
\begin{aligned}
 h(x,y,z)
 &=\frac1{p!}\left[z\left(\frac{\partial}{\partial\lambda_1}\frac{\partial}{\partial x}
 +\frac{\partial}{\partial\lambda_2}\frac{\partial}{\partial y}\right)\right]^p g(\lambda;xy)\\
 &=h_0(xyz)+h_1(xyz)+\cdots+h_p(xyz),
\end{aligned}
\]
donde, por 2.,
\[
  h_i(xyz)=
  \left(z\frac{\partial}{\partial x}\right)^{p-i}
  \left(z\frac{\partial}{\partial y}\right)^i g_i(xy).
\]
Las $h_i$ individuales, y por consiguiente también $h$, son por tanto formas $PZ$ y sumas de ellas; sea $\mathcal T$ de rango $\tau$ respecto de $R$.

Como muestra la construcción de las formas $h(x,y,z)$ de $\mathcal T$, estas formas surgen de $\theta$ mediante procesos polares $P$, y en las filas $x,y,z$ tienen individualmente la misma dimensión que $\theta$; de ahí que la familia $\mathcal T$ sea una subfamilia de $\mathcal L$. Por el hecho mencionado arriba, sólo queda probar la igualdad de los rangos. En esta prueba no se usa ya la estructura especial de los $f(xyz)$, a saber, que son polares de una y la misma forma $\theta$.

\emph{a) Para comparar los rangos de $\mathcal L$ y $\mathcal S$ usamos el Lema a): de $f(x,y,\lambda_1x+\lambda_2y)=0$ se sigue necesariamente que $f(x,y,z)=0$.} Esto se sigue inmediatamente de la relación idéntica entre tres filas binarias
\[
  (xy)z+(yz)x+(zx)y=0,
  \qquad (xy)=(x_1y_2-x_2y_1),\ldots,
\]
que implica
\[
  f[x,y,(xy)x+(xz)y]=(xy)^p f(x,y,z).
\]
Así, de $f(x,y,\lambda_1x+\lambda_2y)=0$ (idénticamente en $\lambda$) se sigue, como muestra la sustitución $\lambda_1=(zy)$, $\lambda_2=(xz)$, y puesto que $(xy)\ne0$, que necesariamente $f(x,y,z)=0$.

Por este lema hay una correspondencia biunívoca entre las formas de $\mathcal S$ y las de $\mathcal L$. Pues de
\[
  f_1(x,y,\lambda_1x+\lambda_2y)=g(\lambda;xy),\qquad
  f_2(x,y,\lambda_1x+\lambda_2y)=g(\lambda;xy)
\]
se sigue necesariamente que
\[
  f_1(xyz)=f_2(xyz).
\]
Además, toda relación en $\mathcal S$ se conserva entre las formas correspondientes de manera única en $\mathcal L$ (y recíprocamente). Pues de
\[
  c_1g^{(1)}(\lambda;xy)+c_2g^{(2)}(\lambda;xy)+\cdots+c_rg^{(r)}(\lambda;xy)=0
\]
se sigue por el lema que necesariamente
\[
  c_1f^{(1)}(x,y,z)+c_2f^{(2)}(x,y,z)+\cdots+c_rf^{(r)}(x,y,z)=0.
\]
\emph{Esto prueba la igualdad de los rangos de $\mathcal L$ y $\mathcal S$; $\rho=\sigma$.}

\emph{b) Para comparar los rangos de $\mathcal S$ y $\mathcal T$ usamos el Lema b) correspondiente: de $h(x,y,z)=0$ se sigue necesariamente que $g(\lambda;xy)=0$.} Para la demostración, obsérvese que, por 3., $h(x,y,z)=0$ es equivalente a la anulación de los productos de potencias individuales
\[
 \left(\frac{\partial}{\partial\lambda_1}\frac{\partial}{\partial x_1}
 +\frac{\partial}{\partial\lambda_2}\frac{\partial}{\partial y_1}\right)^{p_1}
 \left(\frac{\partial}{\partial\lambda_1}\frac{\partial}{\partial x_2}
 +\frac{\partial}{\partial\lambda_2}\frac{\partial}{\partial y_2}\right)^{p_2}g(\lambda;x,y),
 \qquad p_1+p_2=p.
\]
Una diferenciación ulterior da entonces también la anulación de los productos de potencias
\[
\begin{aligned}
&\left(\frac{\partial}{\partial x_1}\right)^{\alpha_1}
 \left(\frac{\partial}{\partial x_2}\right)^{\alpha_2}
 \left(\frac{\partial}{\partial y_1}\right)^{\beta_1}
 \left(\frac{\partial}{\partial y_2}\right)^{\beta_2} \\
&\quad\cdot
 \left(\frac{\partial}{\partial\lambda_1}\frac{\partial}{\partial x_1}
 +\frac{\partial}{\partial\lambda_2}\frac{\partial}{\partial y_1}\right)^{p_1}
 \left(\frac{\partial}{\partial\lambda_1}\frac{\partial}{\partial x_2}
 +\frac{\partial}{\partial\lambda_2}\frac{\partial}{\partial y_2}\right)^{p_2}
 g(\lambda;xy).
\end{aligned}
\]
Por consiguiente, toda combinación lineal de estos productos de potencias se anula, y por tanto toda forma
\[
  f\left(\frac{\partial}{\partial x},\frac{\partial}{\partial y},
  \frac{\partial}{\partial\lambda_1}\frac{\partial}{\partial x}
  +\frac{\partial}{\partial\lambda_2}\frac{\partial}{\partial y}\right)g(\lambda,xy).
\]
Escogiendo $f$ en particular de modo que
\[
  f(x,y,\lambda_1x+\lambda_2y)=g(\lambda;xy),
\]
se obtiene también
\[
  g\left(\frac{\partial}{\partial x};
    \frac{\partial}{\partial x},
    \frac{\partial}{\partial y}\right)g(\lambda,x,y)=0;
\]
o bien, si
\[
  g(\lambda,xy)=\sum G_i\lambda_1^{\nu_1}\lambda_2^{\nu_2}
  x_1^{\alpha_1}x_2^{\alpha_2}y_1^{\beta_1}y_2^{\beta_2},
\]
donde los $G_i$ son formas lineales racionales enteras en las indeterminadas $a$, entonces
\[
  \sum \nu_1!\nu_2!\alpha_1!\alpha_2!\beta_1!\beta_2!\,G_i^2=0,
\]
y por tanto $g(\lambda;xy)=0$.

Del Lema b) se sigue, exactamente como en a), que \emph{los rangos de $\mathcal S$ y $\mathcal T$ son iguales; $\sigma=\tau$.}

Así, por a) y b), se tiene $\rho=\tau$, y por consiguiente ---como $\mathcal T$ es una subfamilia de $\mathcal L$--- queda probada la equivalencia de estas dos familias lineales. Por ello toda forma de $\mathcal L$, en particular también $\theta$, puede representarse como una combinación lineal racional-entera de $\rho$ formas linealmente independientes $h(x,y,z)$:
\[
  \theta=\sum c_i h^{(i)}(x,y,z)=\sum PZ.
\]
Esta identidad, obtenida para coeficientes indeterminados $a$ de $\theta$, sigue siendo válida cuando las indeterminadas $a$ se sustituyen por cantidades de cualquier dominio de racionalidad; por tanto, \emph{el teorema de reducción queda probado en general para el caso de tres filas binarias.}

La demostración general, para $N$ filas de $n$ variables cada una, es completamente paralela. Sea $\theta(A)$ de dimensión $\alpha_i$ en la fila $A_i$. Consideramos de nuevo las tres familias lineales de formas:

1. la familia $\mathcal L$, compuesta por todas las formas $f(A)$ que surgen de $\theta(A)$ mediante procesos polares $P$ y, en cada fila individual $A_i$, tienen la misma dimensión que $\theta(A)$;

2. la familia $\mathcal S$, compuesta por todas las formas $g(\lambda;A_1\cdots A_n)$ que surgen de $f(A)$ mediante las sustituciones
\[
\begin{aligned}
 A_{n+1}&=\lambda_{n+1}^{(1)}A_1+\lambda_{n+1}^{(2)}A_2+
        \cdots+\lambda_{n+1}^{(n)}A_n,\\
 &\vdotswithin{=}\\
 A_N&=\lambda_N^{(1)}A_1+\lambda_N^{(2)}A_2+
        \cdots+\lambda_N^{(n)}A_n;
\end{aligned}
\]
en este proceso, los coeficientes de los productos de potencias individuales de las $\lambda$ en $g(\lambda;A_1\cdots A_n)$ dan formas $Z$ de la identidad (1);

3. la familia $\mathcal T$, compuesta por todas las formas $h(A)$ definidas por
\[
\begin{aligned}
 h(A)=&\left[A_{n+1}\left(
 \frac{\partial}{\partial\lambda_{n+1}^{(1)}}\frac{\partial}{\partial A_1}
 +\cdots+
 \frac{\partial}{\partial\lambda_{n+1}^{(n)}}\frac{\partial}{\partial A_n}\right)\right]^{\alpha_{n+1}} \\
&\cdots
\left[A_N\left(
 \frac{\partial}{\partial\lambda_N^{(1)}}\frac{\partial}{\partial A_1}
 +\cdots+
 \frac{\partial}{\partial\lambda_N^{(n)}}\frac{\partial}{\partial A_n}\right)\right]^{\alpha_N}
 g(\lambda;A_1\cdots A_n).
\end{aligned}
\]
Aquí $h(A)$ es una suma de formas $PZ$.

A causa de la relación idéntica entre cualesquiera $n+1$ filas, el Lema a) sigue siendo válido; igualmente lo sigue siendo el Lema b), ya que el número de variables no entró en su demostración. Así $\mathcal L$ y $\mathcal T$ tienen el mismo rango y ---como $\mathcal T$ es una subfamilia de $\mathcal L$--- son idénticas. En consecuencia, la identidad
\[
  \theta=\sum PZ
\]
vale para coeficientes indeterminados y, por ello, para coeficientes que sean cantidades de cualquier dominio de racionalidad; \emph{con ello queda probado en general el teorema de reducción.}

\subsection*{III.}
Indicamos todavía brevemente cómo consideraciones análogas producen también el \emph{desarrollo en serie general} (para $N\le n$). Sea $Z$ una forma separadamente homogénea en las $n$ filas $A_1,
\ldots,A_n$ de $n$ cantidades cada una; y sea $\Delta$ el determinante de estas filas. Probamos primero la identidad correspondiente a (1):
\begin{equation}
  Z=\sum PH \pmod{\Delta},
\tag{2}
\end{equation}
donde las formas $H$ contienen sólo las filas $A_1,
\ldots,A_{n-1}$ y se derivan de $Z$ mediante procesos $P$.

La demostración consiste en convertir las relaciones obtenidas en II en congruencias módulo $\Delta$. Por simplicidad tomemos tres filas ternarias $x,y,z$; se supone que los coeficientes son indeterminados.

Las tres familias lineales $\mathcal L,\mathcal S,\mathcal T$ se definen como en II para el caso de filas binarias; sean sus rangos $\rho,\sigma,\tau$, mientras que $\rho^*$ y $\tau^*$ denotan los rangos de $\mathcal L$ y $\mathcal T$ módulo $\Delta$ (es decir, hay $\rho^*$, pero no $\rho^*+1$, formas de $\mathcal L$ linealmente independientes mod $\Delta$). Entonces, como en II, se tiene

\emph{Lema a): de $f(x,y,\lambda_1x+\lambda_2y)=0$ se sigue necesariamente que $f(xyz)\equiv0\pmod{\Delta}$ (y recíprocamente).} Pues, por la relación entre cuatro filas ternarias,
\[
  \Delta_1x-\Delta_2y+\Delta_3z=\Delta t,
  \qquad \Delta_1=(yzt),\ldots,
\]
siempre existe una relación lineal módulo $\Delta$ entre tres filas ternarias:
\[
  \Delta_1x-\Delta_2y+\Delta_3z\equiv0\pmod{\Delta},
\]
y por tanto, como en II,
\[
  f(x,y,-\Delta_1x+\Delta_2y)=\Delta_3^p f(x,y,z)\pmod{\Delta}.
\]
Se sigue de $f(x,y,\lambda_1x+\lambda_2y)=0$ (idénticamente en $\lambda$), puesto que $\Delta$ es irreducible y $\Delta_3$ no es divisible por $\Delta$, que necesariamente $f(x,y,z)\equiv0\pmod{\Delta}$. La recíproca es clara a partir de $\Delta(x,y,\lambda_1x+\lambda_2y)=0$. De ahí, como en II, que $\rho^*=\sigma$.

El Lema b) sigue siendo válido con su demostración, y por tanto se tiene $\sigma=\tau$.

Para probar $\tau=\tau^*$ usamos

\emph{Lema c): de $h(x,y,z)\equiv0\pmod{\Delta}$ se sigue necesariamente que $h(x,y,z)=0$.} En efecto, como polar, $h(x,y,z)$ se anula bajo la aplicación del conocido proceso $\Omega$; esto se expresa por la ecuación diferencial
\[
  \Delta\!\left(\frac{\partial}{\partial x},\frac{\partial}{\partial y},\frac{\partial}{\partial z}\right)h(x,y,z)=0.
\]
Si ahora $h(x,y,z)=\Delta(x,y,z)\cdot k(x,y,z)$, entonces por nueva diferenciación se tiene también la anulación de
\[
  k\!\left(\frac{\partial}{\partial x},\frac{\partial}{\partial y},\frac{\partial}{\partial z}\right)
  \Delta\!\left(\frac{\partial}{\partial x},\frac{\partial}{\partial y},\frac{\partial}{\partial z}\right)
  =
  h\!\left(\frac{\partial}{\partial x},\frac{\partial}{\partial y},\frac{\partial}{\partial z}\right)h(x,y,z),
\]
y por tanto de $h(x,y,z)$. Así $\rho^*=\sigma=\tau=\tau^*$; como $\mathcal T$ es una subfamilia de $\mathcal L$, queda probada la identidad (2). El mismo argumento da la demostración para $n$ variables.\footnote{Después de 3. se tiene la relación correspondiente, que se anula a causa del factor determinante, para $\Omega(h)$; en notación operatoria es la ecuación que surge de productos de
\[
 \left(\lambda_1\frac{\partial}{\partial \xi}+\lambda_2\frac{\partial}{\partial \eta}\right),
 \qquad
 \left[z\left(\frac{\partial}{\partial\lambda_1}\frac{\partial}{\partial \xi}
 +\frac{\partial}{\partial\lambda_2}\frac{\partial}{\partial \eta}\right)\right],
\]
y se anula por la anulación del factor determinante.}

Aplicando la identidad (2) al multiplicador de $\Delta$, se obtiene un desarrollo
\[
  Z=\varphi_0+\Delta\varphi_1+\Delta^2\varphi_2+\cdots+\Delta^\mu\varphi_\mu,
\]
donde las $\varphi_i$ se anulan cada una bajo la aplicación del proceso $\Omega$. La repetición $i$ veces del proceso da, puesto que en general $\Omega(\Delta\cdot\psi)=c_1\psi+c_2\Delta\cdot\Omega(\psi)$,
\[
  c\cdot \Omega^i(Z)\equiv\varphi_i\pmod{\Delta};
\]
por otra parte, por (2) se tiene
\[
  c\cdot\Omega^i(Z)\equiv\sum PH_i\pmod{\Delta},
\]
donde los $H_i$ se derivan de la forma $\Omega^i(Z)$ del mismo modo que $H$ se deriva de $Z$. Por el Lema c) (extendido a $n$ variables) se obtiene entonces $\varphi_i=\sum PH_i$, y por consiguiente
\begin{equation}
  Z=\sum PH+\Delta\cdot\sum PH_1+\Delta^2\cdot\sum PH_2+
  \cdots+\Delta^\mu\cdot\sum PH_\mu
\tag{3}
\end{equation}
\emph{como el desarrollo en serie más general de una forma en $n$ filas de $n$ variables cada una, según polares de formas en sólo $n-1$ filas.}

El desarrollo en serie para formas en $N$ filas de $n$ variables $(N<n)$ se obtiene entonces, como es sabido, mediante una sustitución simple. Ponemos (para $N=3$, $n>3$)
\[
  u=\xi_1x+\xi_2y+\xi_3z,
  \quad v=\eta_1x+\eta_2y+\eta_3z,
  \quad w=\zeta_1x+\zeta_2y+\zeta_3z,
\]
y desarrollamos $H(u,v,w)=Z(\xi,\eta,\zeta)$ según (3). Entonces, puesto que $\xi,\eta,\zeta$ sólo aparecen en las combinaciones $u,v,w$, se tiene
\[
  \left(\eta\frac{\partial}{\partial \xi}\right)=
  \left(v\frac{\partial}{\partial u}\right)\cdots,
  \qquad
  \Omega=\sum_{ikl}\left(\frac{\partial}{\partial u}
  \frac{\partial}{\partial v}
  \frac{\partial}{\partial w}\right)_{ikl}(xyz)_{ikl}
  =\nabla_{uvw}(xyz).
\]
Bajo la especialización $\xi_1=\eta_2=\zeta_3=1$; $\xi_2=\xi_3=\cdots=\zeta_2=0$, $H(u,v,w)$ pasa a $H(x,y,z)$, y $\nabla_{uvw}(xyz)$ a $\nabla_{xyz}(xyz)=\nabla_{xyz}$ salvo un factor numérico, ya que la aplicación de $(y\partial/\partial x)$, etc., a los determinantes $(xyz)_{ikl}$ los hace anularse. Puesto que la afirmación correspondiente vale para $N$ filas, (3) da el desarrollo
\begin{equation}
  H=\sum P\Phi+\sum P\Phi_1+\cdots+\sum P\Phi_\mu,
\tag{4}
\end{equation}
donde las $\Phi_i$ surgen de $\nabla^i_{A_1\ldots A_N}H$ mediante procesos $P$ y contienen sólo las filas $A_1,
\ldots,A_{N-1}$, aparte de las combinaciones determinantes que aparecen en $\nabla^i$, las cuales, como se mencionó, no entran en consideración bajo la polarización.

Si esos determinantes en $\nabla^i$ se sustituyen entonces por indeterminadas $p$, las formas resultantes pueden desarrollarse de nuevo según (4); de este modo se obtiene finalmente un desarrollo
\begin{equation}
  H=\left[\sum P\psi(p)\right]_{p=\text{determinantes de }A_1\cdots A_N},
\tag{5}
\end{equation}
donde las $\psi$ surgen de la forma original mediante los procesos $P$ y $\nabla_{A_1\ldots A_N}(p)$, $\nabla_{A_1\ldots A_{N-1}}(p)$, etc. Con estos desarrollos y sus combinaciones ---por ejemplo, aplicando (3), y luego (4) y (5), a las formas $Z$ que aparecen en (1)--- se agotan todos los desarrollos conocidos para formas en arbitrariamente muchas filas de $n$ variables cogredientes cada una.

Erlangen, 5 de enero de 1915.

\end{document}
